\chapter{Contribuția proprie}
\label{cap:contributie}

\section{Cerințe și prezentare generală a sistemului}

Pe scurt, ce poate face un utilizator autentificat: își gestionează datele de planificare (un catalog
de exerciții oficiale, comun, plus exerciții proprii, rutine de încălzire/încheiere, săli și
echipamente), își construiește \emph{programe} de antrenament pe mai multe zile, execută sesiuni live
(unde notează seriile cu greutate, repetări, RPE), își vede istoricul și progresul, publică și copiază
programe într-un \emph{marketplace}, primește o analiză a structurii unui program și, dacă este
antrenor, poate vedea (doar citi) datele unui client cu care are o relație activă. Peste programe și
peste istoric funcționează o \emph{căutare} full-text.

Cerințele nefuncționale au ghidat tot ce am făcut: PostgreSQL este sursa unică de adevăr; schema este
fixată de la început și se schimbă doar prin migrări Flyway; fiecare resursă a unui utilizator
verifică proprietarul, ca să nu existe IDOR; iar codul este gândit la o calitate de tip „portofoliu /
sisteme back-end”, fără substitute în memorie, fără H2, cu teste de integrare reale. Câteva capturi
din interfață se găsesc în \textbf{Anexa~1}.

\section{Arhitectura sistemului}
\label{sec:arhitectura}

Sistemul este o aplicație web pe trei niveluri: o aplicație React (single-page) vorbește prin JSON cu
un API REST scris în Spring Boot, care salvează în PostgreSQL. În producție, frontend-ul compilat este
inclus în backend și servit de la aceeași adresă (\code{http://localhost:8080/}), așa că tot cookie-ul
de autentificare rămâne same-origin. Tot stack-ul rulează prin Docker Compose (\textbf{Anexa~6}).

Backend-ul este organizat \emph{pe funcționalități}, nu pe niveluri tehnice. Fiecare pachet
(\code{auth}, \code{exercise}, \code{routine}, \code{gym}, \code{template}, \code{session},
\code{history}, \code{analytics}, \code{marketplace}, \code{analyzer}, \code{search}, \code{coaching},
\code{shared}) își are propriile patru niveluri: \code{domain} (entitățile JPA), \code{application}
(servicii, comenzi, excepții), \code{infrastructure} (depozite Spring Data și SQL nativ) și \code{web}
(controllere și DTO-uri). O regulă pe care am ținut-o strict: nivelul web nu expune niciodată
entități, ci doar DTO-uri. Când o funcționalitate are nevoie de datele alteia (de exemplu analizorul
citește structura unui program), cel care le deține expune un \emph{model de citire la nivel de
aplicație}, nu un DTO web și nici un depozit; așa dependențele rămân curate.

O cerere este autentificată de un filtru care verifică JWT-ul și pune în contextul de securitate un
\code{UserPrincipal}. Controllerele citesc acest principal și nu se încred niciodată într-un
identificator de utilizator venit din corpul cererii. Erorile de domeniu extind o excepție comună
care duce cu ea un cod HTTP și un cod de eroare stabil, ușor de citit de un program, iar un singur
handler le transformă într-un format de eroare uniform (\textbf{Anexa~5}). Diagrama de arhitectură
este în Figura~\ref{fig:arhitectura}.

\begin{figure}[!htb]
  \capturaplaceholder{Diagrama de arhitectură: SPA React $\rightarrow$ API REST Spring Boot
  $\rightarrow$ PostgreSQL (sursa de adevăr) și OpenSearch (model de citire separat).}
  \caption{Arhitectura de ansamblu a sistemului.}
  \label{fig:arhitectura}
\end{figure}

\section{Modelul de date și principiul instantaneelor}
\label{sec:instantanee}

Schema este creată de o singură migrare inițială (\code{V1}), care acoperă utilizatorii, catalogul de
exerciții și grupele musculare, rutinele, sălile și echipamentele, programele cu tot ce conțin
(zile / exerciții / rutine), sesiunile live cu instantaneele lor, tabelele de marketplace (voturi,
salvări, evenimente de utilizare, statistici per program) și câteva tabele lăsate pentru mai târziu
(coaching, rezultate de analiză, un \emph{outbox}). A doua migrare (\code{V2}) pune catalogul oficial
de exerciții, iar a treia (\code{V3}) adaugă un index pentru concurență. Toate tabelele și
constrângerile sunt în \textbf{Anexa~2}.

\paragraph{Principiul instantaneelor.} Faptul că istoricul rămâne corect este cerința mea de bază și o
asigur în două locuri. Întâi, exercițiile dintr-o zi de program rețin
\code{exercise\_name\_snapshot}, \code{exercise\_type\_snapshot} și o copie a grupelor musculare,
păstrând doar o legătură \emph{care poate fi anulată} către exercițiul viu (\code{ON DELETE SET
NULL}). Apoi, când pornește un antrenament, fac un al doilea instantaneu: numele programului, al zilei
și al sălii, plus numele/tipul/valorile planificate ale fiecărui exercițiu și grupele musculare, sunt
copiate în \code{session\_exercises}; seriile rețin și numele echipamentului. Pentru că rândurile
sesiunii își poartă propriile valori și au doar chei externe care pot fi anulate, dacă editez sau
șterg programul, exercițiul, rutina, sala sau echipamentul, ce a fost notat într-o sesiune veche nu se
schimbă. Asta o verific direct prin teste.

\paragraph{Ștergere logică, unicitate și constrângeri.} Resursele de planificare folosesc o coloană
\code{deleted\_at}, ca istoricul care le referă să supraviețuiască. Unicitatea o impun cu \emph{indecși
unici parțiali} care ignoră rândurile șterse (de exemplu, un nume unic de exercițiu propriu per
utilizator \code{WHERE deleted\_at IS NULL}), așa că un nume poate fi reutilizat după ștergere. Multe
reguli sunt impuse chiar de baza de date prin constrângeri \code{CHECK}: valori permise pentru
enumerări, \code{days\_per\_week BETWEEN 1 AND 7}, \code{rpe BETWEEN 1.0 AND 10.0},
\code{finished\_at >= started\_at}, contoare care nu pot fi negative, interdicția de a te antrena
singur (\code{coach\_user\_id <> client\_user\_id}) și regula că un program public trebuie să aibă o
dată de publicare.

\paragraph{Concurență.} Regula „un singur antrenament activ per utilizator” este garantată de un
\emph{index unic parțial} \code{ux\_workout\_sessions\_one\_active\_per\_user … WHERE status =
'IN\_PROGRESS'} (migrarea \code{V3}). Verificarea din serviciu este doar o cale rapidă și prietenoasă;
dacă două cereri pornesc în același timp, cea care pierde primește un \code{409 ACTIVE\_WORKOUT\_EXISTS}
curat, nu o eroare 500. Contoarele de marketplace (voturi, salvări) sunt actualizate cu SQL atomic și
mărginit (\code{GREATEST(0, contor +/- delta)}), deci nu apare cursa citire–modificare–scriere și nu
ajung niciodată sub zero.

\section{Modelul de securitate}
\label{sec:securitate}

Parolele sunt stocate ca amprente BCrypt. La autentificare, API-ul dă un token de acces JWT de scurtă
durată (HS256), pe care SPA-ul îl trimite în antetul \code{Authorization} și îl ține doar în memorie,
și un token de reîmprospătare opac, ca un cookie \code{HttpOnly}, \code{SameSite=Strict}, limitat la
calea \code{/api/auth}. Din token-ul de reîmprospătare păstrez în baza de date doar o amprentă SHA-256.
La reîmprospătare, token-ul este \emph{rotit}: cel prezentat este verificat și revocat, iar o pereche
nouă este emisă, totul într-o singură tranzacție. Așa, dacă cineva fură un token și încearcă să-l
refolosească, este prins (amprenta lui există, dar are \code{revoked\_at} setat). Am dezactivat
protecția CSRF intenționat, fiindcă token-ul de acces vine din antet, nu din cookie; singurul cookie cu
credențiale este cel de reîmprospătare, protejat de atributele de mai sus.

\paragraph{Autorizare și siguranță IDOR.} Fiecare resursă a unui utilizator este încărcată cu o
interogare care include și identificatorul proprietarului. Important, dacă cineva cere o resursă a
altui utilizator, primește \textbf{404, nu 403}, ca să nu afle nici măcar că resursa există. Resursele
imbricate își verifică proprietarul în lanț (de exemplu, un exercițiu dintr-o zi de program este
accesibil doar dacă programul acelei zile îți aparține). În configurația de securitate, ordinea este:
rute publice de autentificare; \code{/api/admin/**} doar pentru rolul \code{ADMIN};
\code{/api/coach/**} doar pentru rolul \code{COACH}; restul \code{/api/**} doar pentru utilizatori
autentificați.

\section{Funcționalitățile aplicației}
\label{sec:functionalitati}

\paragraph{Date de planificare și programe.} Utilizatorii își fac o bibliotecă de exerciții (catalogul
oficial plus exerciții proprii, fiecare cu grupe musculare primare/secundare), rutine, săli și
echipamente. Un program este privat și pe mai multe zile; fiecare zi are exerciții ordonate (cu
serii/repetări/greutate/pauză planificate) și rutine atașate. După fiecare modificare a exercițiilor,
recalculez pe partea de PostgreSQL niște vectori agregați ai programului (grupe musculare,
identificatori de exerciții oficiale, nume de exerciții), printr-o singură instrucțiune nativă cu
\code{array\_agg(DISTINCT … ORDER BY …)}, ca să pregătesc căutarea.

\paragraph{Execuția live și istoricul.} Pornirea unui antrenament face, într-o singură tranzacție,
copierea-instantaneu din Secțiunea~\ref{sec:instantanee}. În timpul sesiunii, utilizatorul notează
serii, adaugă exerciții în plus, neplanificate, și la final finalizează sau anulează sesiunea.
Istoricul listează sesiunile cele mai recente primele, cu câteva numere de sumar per sesiune calculate
printr-o singură interogare în lot (ca să evit N+1).

\paragraph{Analize de progres.} Analizele sunt calculate complet în PostgreSQL, peste sesiunile
finalizate: volumul total pe zi, numărul de antrenamente pe săptămână (\code{date\_trunc}), distribuția
seriilor pe grupa musculară primară și, pentru un grafic al progresului la forță, cea mai bună estimare
a unei repetări maxime (1RM) pe exercițiu pe zi, cu formula lui Epley \cite{epley1985}:
\begin{equation}
  \mathrm{1RM} \approx w \cdot \left(1 + \frac{r}{30}\right),
  \label{eq:epley}
\end{equation}
unde $w$ este greutatea și $r$ numărul de repetări ale seriei.

\paragraph{Marketplace.} Proprietarii publică un program nevid, făcându-l public. Utilizatorii
autentificați parcurg programele publice (sortate după cele mai noi, scor sau tendință), votează,
salvează și \emph{folosesc} un program. „Folosirea” copiază în profunzime programul public într-un
program nou, privat, al celui care copiază. Legătura către exercițiu o păstrez doar dacă este un
exercițiu \emph{oficial}; altfel o anulez, dar păstrez copia numelui/tipului/grupelor musculare — așa
copia rămâne complet utilizabilă și independentă de original, fără să scape date private.

\section{Analizorul structural informat de dovezi}
\label{sec:analizor}

Analizorul este determinist și bazat pe reguli — nu folosește învățare automată. Dă un \emph{Scor de
structură a programului} (nu o judecată de „calitate științifică” sau medicală), iar fiecare răspuns
vine cu un avertisment că nu este sfat medical și că nu ține cont de recuperare, tehnică, efort, somn,
alimentație sau de cum răspunde fiecare individ. Îl calculez la cerere, la un \code{GET} read-only
(\code{/api/templates/\{id\}/analysis}), și îl folosesc și la indexarea pentru căutare.

Volumul săptămânal este suma peste zilele scrise în program, iar seriile sunt ponderate după rolul
muscular (primar 1.0, secundar 0.5). Exercițiile fără serii planificate sunt \emph{scoase} din volum
(nu le presupun o valoare implicită) și marcate ca date incomplete; doar exercițiile de forță contează
la volum și echilibru. Scorul de 100 de puncte este împărțit în cinci părți, fiecare pornind de la
maximul ei și scăzând o valoare fixă pentru fiecare problemă găsită (dar nu sub zero), deci rezultatele
sunt complet deterministe și pot fi testate (Tabelul~\ref{tab:analizor}). Categoriile pun totalul în
\code{WELL\_STRUCTURED} ($\geq 80$), \code{DECENT\_STRUCTURE} ($55$–$79$) sau \code{NEEDS\_REVIEW}
($<55$). Familiile de reguli și o bucată reprezentativă din logica de scor sunt în \textbf{Anexa~3}.

\begin{table}[!htb]
  \centering
  \small
  \begin{tabular}{@{}lcl@{}}
  \toprule
  Sub-scor & Maxim & Exemple de penalizări \\
  \midrule
  Volum și acoperire & 35 & regiune absentă $-12$; volum excesiv $-5$; foarte scăzut $-3$ \\
  Frecvență & 20 & volum concentrat într-o singură zi $-4$ \\
  Echilibru & 20 & dezechilibru sever $-8$; moderat $-4$ \\
  Designul sesiunii & 15 & sesiune excesivă $-5$; concentrare pe o grupă $-3$ \\
  Specificitate / pauză & 10 & pauză prea scurtă $-2$; interval de repetări extrem $-2$ … $-3$ \\
  \bottomrule
  \end{tabular}
  \caption{Descompunerea scorului de structură pe cinci sub-scoruri deterministe.}
  \label{tab:analizor}
\end{table}

\section{Modelul de citire pentru căutare (OpenSearch)}
\label{sec:cautare}

Căutarea este un \emph{model de citire separat}, construit peste OpenSearch și pornit printr-un
indicator de configurare (\code{app.search.enabled}, implicit oprit, ca să nu lege compilarea locală
și testele pe PostgreSQL de OpenSearch). Am două indexuri, \code{templates} și \code{workout\_sessions},
cu nume versionate ascunse în spatele unor alias-uri stabile. Analizorul de indexare normalizează
textul (lowercase + scoaterea diacriticelor), iar cel de căutare adaugă și un set de \emph{sinonime}
(de exemplu \emph{pecs}$\rightarrow$\emph{chest}, \emph{ohp}$\rightarrow$\emph{overhead press}).
Interogarea de tip \code{multi\_match} dă ponderi (boost) câmpurilor — pentru programe \code{name\^{}4},
\code{exerciseNames\^{}3}, \code{muscleGroupsText\^{}2} etc. — și toleranță la greșeli
(\code{fuzziness=AUTO}, cu o regulă care nu aplică „fuzzy” pentru căutări mai scurte de 3 caractere).
Pentru marketplace, scorul de relevanță este combinat cu un \emph{function\_score} care adaugă o
influență logaritmică a popularității (salvări, utilizări). Rezultatele includ fațete (agregări
\code{terms} și un histogram pe lună) și evidențierea potrivirilor cu \code{<mark>}.

\paragraph{Cum scriu în index și cum țin totul consistent.} Indexarea este pornită de evenimente
declanșate \emph{după} ce tranzacția se salvează (\code{@TransactionalEventListener(AFTER\_COMMIT)}),
ca să nu ajungă în index munca anulată prin rollback. Am două căi de actualizare: o reindexare
\emph{structurală} reconstruiește tot documentul și recalculează câmpurile din analiză, iar o
actualizare \emph{doar a contoarelor} (vot/salvare/utilizare) nu mai recalculează analiza. Toate
scrierile sunt \emph{idempotente} (rescriu documentul întreg după identificator). Dacă o scriere dă
greș, o jurnalizez și merg mai departe — PostgreSQL rămâne sursa de adevăr, iar o reconstrucție repară
indexul. Intenționat, aceasta \emph{nu} este încă un outbox tranzacțional (vezi
Secțiunea~\ref{sec:directii}).

\paragraph{Securitate pe două straturi.} Întâi, un \emph{filtru obligatoriu}, luat din JWT, pe care
parametrii cererii nu îl pot suprascrie: \code{marketplace} $\Rightarrow$ \code{visibility=PUBLIC};
\code{my} $\Rightarrow$ \code{ownerUserId=apelant}; \code{workouts} $\Rightarrow$
\code{ownerUserId=apelant}. Apoi, o verificare suplimentară în PostgreSQL: identificatorii întorși de
OpenSearch sunt verificați din nou în baza de date cu aceeași regulă, ca să arunc orice rezultat
învechit (de exemplu, abia depublicat sau șters) și să adaug starea curentă de vot/salvare pentru
marketplace.

\paragraph{Reconstrucție și reziliență.} Serviciul de reconstrucție creează un index nou cu marcaj de
timp, îl umple în lot din PostgreSQL și \emph{schimbă atomic alias-ul} (fără să se oprească citirile),
apoi șterge indexul vechi; dacă ceva pică, indexul nou este aruncat, iar cel vechi continuă să
servească. Reconstrucția manuală (\code{POST /api/admin/search/reindex}) este permisă doar rolului
\code{ADMIN}. Dacă căutarea este oprită sau OpenSearch nu răspunde, endpoint-urile întorc \code{503} și
interfața trece înapoi pe lista din SQL; un index încă neconstruit este tratat ca „fără rezultate”, nu
ca o eroare.

\paragraph{Căutarea făcută de antrenor.} Un antrenor poate căuta full-text în istoricul unui client
(\code{/api/coach/clients/\{id\}/search/workouts}). Este singurul loc unde căutarea rulează cu un
identificator de proprietar diferit de cel care cere, și este sigură fiindcă trece mai întâi prin
poarta de acces a antrenorului (vezi Secțiunea~\ref{sec:coaching}) și apoi folosește exact același
serviciu de căutare, dar cu identificatorul clientului, păstrând filtrul obligatoriu și verificarea din
PostgreSQL.

\section{Evaluarea empirică: OpenSearch vs.\ PostgreSQL}
\label{sec:benchmark}

Partea cea mai importantă din punct de vedere științific este o măsurătoare concretă: de la ce
dimensiune merită un motor de căutare dedicat față de ce poate face PostgreSQL singur.

\paragraph{Metodologie.} Varianta de PostgreSQL este făcută \emph{corect}, nu ca să piardă: căutare
full-text ponderată cu \code{tsvector} (\code{setweight} A/B/C/D) și index GIN, ordonată cu
\code{ts\_rank\_cd}; căutare fuzzy cu \code{pg\_trgm} și index GIN; fațete prin \code{GROUP BY}. Partea
OpenSearch folosește exact aceleași mapări și interogări ca în producție. Setul de date este generat
determinist și \emph{în flux} (pe loturi, cu memorie constantă), cu un vocabular realist de ridicări,
popularitate înclinată și 90\%/10\% public/privat. Am măsurat la șapte dimensiuni, de la 1.000 la
2.000.000 de documente, pe cinci scenarii: două căutări full-text, o căutare fuzzy (greșeală de
tastare), o căutare filtrată și o agregare pe fațete. Măsurarea este pe un singur fir, cu încălzire și
un buget de timp, și raportez percentilele $p_{50}/p_{95}/p_{99}$.

\paragraph{Rezultate.} Tabelul~\ref{tab:bench} arată latențele la cea mai mare dimensiune (2M de
documente); tabelele complete, pe toate dimensiunile, sunt în \textbf{Anexa~4}. Pe scurt: cele două
soluții se echivalează în jurul a \emph{10.000 de documente} (sub această valoare PostgreSQL e mai
rapid; peste ea, OpenSearch câștigă clar); căutarea full-text în PostgreSQL crește aproape liniar (de
la $\sim$1,5\,ms la $\sim$1,86\,s între 1k și 2M), în timp ce OpenSearch rămâne aproape constant
($\sim$3–64\,ms); \emph{căutarea fuzzy face diferența} — \code{pg\_trgm} se duce de la 14\,ms la
$\sim$15,8\,s (1k$\rightarrow$2M), față de $\sim$4\,ms, aproape constant, în OpenSearch (un raport de
ordinul a 4000$\times$ la 2M); în plus, indexul OpenSearch ocupă de $\sim$3,6$\times$ mai puțin pe
disc și se construiește de $\sim$2$\times$ mai repede la 2M.

\begin{table}[!htb]
  \centering
  \begin{tabular}{@{}lrr@{}}
  \toprule
  Scenariu (la 2M documente) & PostgreSQL & OpenSearch \\
  \midrule
  Full-text „bench press” ($p_{50}$) & 1864\,ms & 64\,ms \\
  Full-text „squat” ($p_{50}$) & 1437\,ms & 4\,ms \\
  Fuzzy „benhc” (greșeală) ($p_{50}$) & 15825\,ms & 4\,ms \\
  Filtrat (press + INTERMEDIATE + PPL) ($p_{50}$) & 2295\,ms & 7\,ms \\
  Fațetă după dificultate ($p_{50}$) & 1193\,ms & 2\,ms \\
  \midrule
  Timp de construcție a indexului & 239\,s & 115\,s \\
  Dimensiune pe disc & 2632\,MB & 739\,MB \\
  \bottomrule
  \end{tabular}
  \caption{Comparație OpenSearch vs.\ PostgreSQL la 2.000.000 de documente. Tabelele complete pe
  toate cele șapte dimensiuni sunt în Anexa~4.}
  \label{tab:bench}
\end{table}

\paragraph{Interpretare.} Rezultatul transformă o alegere de arhitectură luată des „din reflex” într-o
decizie cu temei: sub pragul de echivalență, complexitatea în plus a unui motor dedicat nu se justifică;
peste el — și mai ales dacă ai nevoie de toleranță la greșeli — un motor dedicat devine nu doar util, ci
necesar.

\section{Modul de coaching: RBAC plus ReBAC}
\label{sec:coaching}

Modul de coaching adaugă un al doilea strat de control peste RBAC. O relație antrenor–client are un
ciclu de viață (\code{PENDING} $\rightarrow$ \code{ACTIVE} dacă se acceptă, sau \code{REJECTED}; o
relație \code{ACTIVE} poate fi \code{REVOKED} de oricare parte), cu un index unic parțial care permite
o singură relație vie per pereche. RBAC (\code{/api/coach/**} cere rolul \code{COACH}) este necesar,
dar nu suficient: fiecare citire per client trece printr-o poartă bazată pe relații (ReBAC) care cere o
relație \emph{activă} și, altfel, întoarce \textbf{404 (nu 403)}, ca un antrenor să nu poată afla ce
utilizatori există. Accesul antrenorului este doar de citire (istoric, analize, o sesiune, căutare); nu
există niciun endpoint prin care antrenorul să modifice datele clientului, iar clientul poate retrage
accesul oricând.

\section{Strategia de testare}
\label{sec:testare}

Corectitudinea o verific cu teste de integrare pe \emph{infrastructură reală}: o clasă de bază
pornește un container PostgreSQL (Testcontainers), iar suita de căutare adaugă un container OpenSearch
real și pornește indicatorul de căutare — fără mock-uri, fără H2. Așa, toată schema gestionată de
Flyway și constrângerile ei sunt puse la treabă. Testele verifică comportamentul prin nivelul HTTP și,
unde contează, garanțiile la nivel de bază de date: o a doua inserare „în desfășurare” brută este
respinsă de indexul unic parțial; vectorii agregați sunt citiți înapoi cu \code{JdbcTemplate};
instantaneele unei sesiuni rămân neschimbate după ce editez datele sursă. La momentul scrierii, suita
are 128 de teste care trec.

\section{Direcții de dezvoltare proiectate}
\label{sec:directii}

Lucrurile de mai jos sunt \emph{gândite și pregătite}, dar încă neimplementate. Le includ aici fiindcă
arhitectura și, în mai multe cazuri, chiar schema bazei de date le anticipează: schema \code{V1}
conține deja, nefolosite de cod, tabelele \code{outbox\_events}, \code{template\_analysis\_results},
\code{coach\_session\_comments} și \code{coach\_template\_assignments}.

\begin{itemize}
  \item \textbf{Cache cu Redis pentru marketplace (cache-aside).} Aș memora paginile de listă și
        detaliile de program, cu invalidare condusă de aceleași evenimente după salvare deja existente,
        plus protecție împotriva efectului de „turmă” prin expirare probabilistică timpurie
        (XFetch \cite{vattani2015stampede}).
  \item \textbf{Outbox tranzacțional pentru indexare.} Ar rezolva problema \emph{dual-write}: un rând
        de eveniment scris în aceeași tranzacție cu modificarea, în tabela \code{outbox\_events} deja
        existentă, plus un proces care îl trimite mai departe cu livrare cel-puțin-o-dată și un
        consumator idempotent (indexatorul este deja idempotent). Aș discuta și alternativa
        CDC/Debezium \cite{richardson2018microservices}.
  \item \textbf{Evaluarea calității relevanței.} Aș duce evaluarea de la viteză la calitate: un set de
        judecăți etichetate și indicatori ca precision@k și nDCG \cite{jarvelin2002ndcg}, plus un
        studiu de ablație (sinonime / fuzzy / ponderi pornite sau oprite).
  \item \textbf{Paginare prin chei (keyset).} Aș înlocui paginarea \code{OFFSET} de acum cu paginare
        prin cheie de continuare, ducând costul paginilor adânci de la $\bigO{n}$ la $\bigO{\log n}$.
  \item \textbf{Limitarea ratei la autentificare} (token-bucket cu Redis) și \textbf{trecerea la
        Argon2id} pentru parole, cu re-hashing transparent la autentificare.
  \item \textbf{Observabilitate} (trasare distribuită OpenTelemetry + metrici Prometheus/Grafana),
        \textbf{partiționarea pe timp} a tabelei \code{workout\_sessions}, o \textbf{coadă de lucru}
        bazată pe \code{SELECT … FOR UPDATE SKIP LOCKED} (potrivită pentru trimiterea din outbox) și
        \textbf{vederi materializate} pentru analize.
\end{itemize}

Aceste direcții sunt descrise, ca proiectare, în Capitolul~\ref{cap:concluzii}.
